\subsection{Conductors over more general local bases}

The residue field and the regularity of the conductor play different
roles in the tensor obstruction.  The next two examples separate them.

\begin{proposition}\label{prop:infinite-residue-nilpotent-conductor}
There is a complete one-dimensional Noetherian Cohen--Macaulay local
ring $A$ and an inclusion $A\subset C$ with a selected nonreduced
Artinian factor and a nonzero conductor, consisting entirely of
zero divisors, for which every polynomial tensor map is an
isomorphism.  Its residue field is infinite.
\end{proposition}
\begin{proof}
Let $k$ be an infinite field, $D=k[[t]]$, and $K=k((t))$.  Set
\[
 A=D[\varepsilon,\eta]/(\varepsilon^2,\varepsilon\eta,\eta^2),
 \qquad
 C=(K[\varepsilon]/\varepsilon^2)
       \times(D[\eta]/\eta^2).
\]
Embed $A$ in $C$ by
\[
 a+b\varepsilon+c\eta\longmapsto
             (a+b\varepsilon,a+c\eta).
\]
This is an injective unital ring homomorphism.  The ring $A$ is
finite free over $D$, complete local with maximal ideal
$(t,\varepsilon,\eta)$, and has reduced quotient $D$.
Thus it is Noetherian of dimension one.  The parameter $t$ is
regular, proving the Cohen--Macaulay assertion.

For $e=(1,0)$, the selected image is
$A_0=eA=D[\varepsilon]/\varepsilon^2$ and the selected factor is
$Q_0=eC=K[\varepsilon]/\varepsilon^2=Q(A_0)$.
Take $O=A_0$.  Its conductor into $A$ is
\[
 \{x\in O:exO\subset A\}=D\varepsilon.
\]
Indeed, if $(a+b\varepsilon,0)$ belongs to $A$, its common scalar
coefficient is zero, so $a=0$; conversely every supported
$\varepsilon$-multiple lies in $A$.  The conductor is nonzero,
square zero, and contains no regular parameter multiple of $O$.

For any polynomial of degree at most $d$ in an arbitrary coefficient
extension of $A$, choose $d+1$ distinct elements of $k$.  Their
Vandermonde determinant is a unit of $A$.  If all polynomial values
belong to $A$, the inverse Vandermonde matrix therefore puts every
coefficient in $A$.  Tensor products of these matrices prove the
same assertion for a finite grid in several variables.  Hence
\[
 \Omega_n(A,C)=A[X_1,\ldots,X_n]
\]
for every $n$, and every tensor map is an isomorphism.
\end{proof}

\begin{proposition}\label{prop:nonnoetherian-nilpotent-conductor}
Let $k=\mathbb F_q$, $D=k[[t]]$, $K=k((t))$, and let $N$ be a
$D$-module.  Give
\[
 A=D\oplus D\varepsilon\oplus N,
 \qquad A_1=D\oplus N
\]
the indicated square-zero extensions of $D$, and embed $A$ in
\[
 C=(K[\varepsilon]/\varepsilon^2)\times A_1,
 \qquad
 d+a\varepsilon+n\longmapsto(d+a\varepsilon,d+n).
\]
The binary tensor map for $\Omega(A,C)$ is not injective.
For $N=\prod_{j\ge1}D$, the ring $A$ is complete separated local
and non-Noetherian, with finite residue field and a nonzero
square-zero selected conductor.
\end{proposition}
\begin{proof}
The map is a unital injection, and the selected order for $e=(1,0)$
is $A_0=D[\varepsilon]/\varepsilon^2$.  Its conductor is
$D\varepsilon$, by the same common-scalar calculation as above.

Write the first component of $F\in\Omega(A,C)$ as
$f_0+\varepsilon f_1$, with $f_i\in K[X]$.
The scalar component of its second polynomial lies in $D[X]$.
Equality of the two scalar values on the infinite domain $D$
forces $f_0\in D[X]$.  Evaluation at $d\in D$ then gives
$f_1(d)\in D$ for every $d$, so $f_1\in\operatorname{Int}(D)$.
Define
\[
 \lambda:\Omega(A,C)\longrightarrow k,
 \qquad \lambda(F)=f_1(t)-f_1(0)\pmod {tD}.
\]
This is $A$-linear through $A\to k$.  In fact, multiplication by
$d+a\varepsilon+n$ changes $f_1$ to $df_1+af_0$, and
$f_0(t)-f_0(0)\in tD$.

The polynomial
\[
 H=\left(\frac{\varepsilon(X^q-X)}t,0\right)
\]
belongs to $\Omega(A,C)$ because $d^q-d\in tD$ for $d\in D$
and the nilpotent part of the argument is annihilated by
$\varepsilon$.  Its value under $\lambda$ is
$t^{q-1}-1\pmod t=-1$.
Since $H(X)H(Y)=0$, the tensor $H\otimes H$ is in the binary
kernel.  It is nonzero: the $A$-balanced map
\[
 \Omega(A,C)\otimes_A\Omega(A,C)\longrightarrow k,
 \qquad F\otimes G\longmapsto\lambda(F)\lambda(G)
\]
sends it to $1$.

For the specified product module, $A$ is local and
\[
 \mathfrak m^r=t^rD\oplus t^{r-1}(D\varepsilon\oplus N)
 \qquad(r\ge1).
\]
Both $D$ and $\prod_{j\ge1}D$ are complete and separated
$t$-adically; the displayed equality proves the same assertion
for the maximal-ideal topology of $A$.
The vector space $N/tN=\prod_{j\ge1}k$ is infinite-dimensional.
Thus $N$ is not finitely generated as a $D$-module.  The
square-zero ideal $N\subset A$ has its $A$-action through $D$,
so it is not finitely generated as an $A$-ideal either.
Consequently $A$ is not Noetherian.
\end{proof}
